% Context from chapters/machine-learning.tex; for mathematical reference, not compilation.
\section{Unsupervised Learning}

In unsupervised learning, we observe $n$ points $(x_i)_{i=1}^n$.
 
The aim is to identify structure in these observations, for example to visualize them or group them into clusters.
 
We assume for simplicity that the observations lie in Euclidean space, $x_i \in \RR^p$, where $p$ is the number of features. We store the observations as the rows of a matrix $X\in\RR^{n\times p}$.


                                                  
\subsection{Dimensionality Reduction and PCA}

Dimensionality reduction produces a compact representation of the data for visualization or subsequent analysis. By retaining informative features, it can also reduce computational cost and improve prediction.
 
Principal component analysis (PCA) projects the data orthogonally onto the principal axes of the covariance matrix.

   
\paragraph{Presentation of the method.}

The empirical mean is defined as 
\eq{
    \hat m \eqdef \frac{1}{n} \sum_{i=1}^n x_i \in \RR^p
} 
and the empirical covariance as
\eql{\label{eq-emp-cov}
 	\hat C \eqdef \frac{1}{n} \sum_{i=1}^n (x_i-\hat m) (x_i-\hat m)^* \in \RR^{p \times p}. 
}
With the centered data matrix $\tilde X \eqdef X-\ones_n\hat m^*$, the covariance is $\hat C=\tilde X^*
\tilde X/n$. 

Suppose the points $(x_i)_i$ are i.i.d. with a finite second moment, and let $\xp$ denote a random variable with their common distribution. The law of large numbers gives the following almost sure limits as $n \rightarrow +\infty$:
\eql{\label{eq-cov-approx}
	\hat m \rightarrow m \eqdef \EE(\xp)
	\qandq
	\hat C \rightarrow C \eqdef \EE((\xp-m)(\xp-m)^*).
}
Let $\mu$ be the distribution on $\RR^p$ of $\xp$. Then these limits can also be written as
\eq{
	m = \int_{\RR^p} x \d\mu(x)
	\qandq
	C = \int_{\RR^p} (x-m)(x-m)^* \d\mu(x).	
}


\begin{figure}
\centering
\begin{tabular}{@{}c@{}c@{}}
\includegraphics[width=.25\linewidth]{ml/pca-nn/cov}&
\includegraphics[width=.4\linewidth]{ml/pca-nn/svd}\\
$C$ & SVD of $C$
\end{tabular}
\caption{\label{fig-cov}
Empirical covariance of the data and its associated singular values. 
}
\end{figure}


\begin{figure}[tbp]
\centering
\BookDrawing{tikz/ml-pca-variance}
\caption{Principal directions of variation.}
\label{fig-review-ml-pca-variance}
\end{figure}
The PCA basis consists of the right singular vectors of the centered data matrix. Proposition~\ref{prop-svd} gives the reduced singular value decomposition
\eq{
 	\tilde X = \sqrt{n} U \diag(\si) V^*
}
where $U \in \RR^{n \times r}$ and $V \in \RR^{p \times r}$ have orthonormal columns, and $r=\rank(\tilde X) \leq \min(n,p)$.
 
Write $V=(v_k)_{k=1}^r$ for the orthonormal columns, which are eigenvectors of $\hat C = V \diag(\si^2) V^\top$, with $v_k \in \RR^p$. These vectors describe the principal directions of variation in the point cloud $(x_i)_i$ in $\RR^p$.
 
We order the singular values as $\si_1 \geq \ldots \geq \si_r$, so the leading components account for the greatest variance.

Figure~\ref{fig-cov} displays an empirical covariance matrix and its spectrum $(\si_k^2)_k$ for the Iris dataset\footnote{\url{https://en.wikipedia.org/wiki/Iris_flower_data_set}} introduced by Fisher. It contains 50 samples from each of three iris species: Iris setosa, Iris virginica, and Iris versicolor. The four features are sepal length, sepal width, petal length, and petal width.

To obtain a PCA embedding $x_i \in \RR^p \mapsto z_i \in \RR^d$ of dimension $d \leq p$, project onto the first $d$ right singular vectors, extending the basis if $d>r$:
\eql{\label{eq-pca-1}
	z_i \eqdef ( \dotp{x_i-\hat m}{v_k} )_{k=1}^d \in \RR^d.
}
The corresponding reconstruction is
\eql{\label{eq-pca-2}
	\tilde x_i\eqdef\hat m+\sum_{k=1}^d z_{i,k}v_k \in \RR^p.
}
Thus $\tilde x_i = \Proj_{\tilde T}(x_i)$, where $\tilde T\eqdef\hat m+ \Span_{k=1}^d(v_k)$ is the affine reconstruction space.

Figure~\ref{fig-pca} shows an example of PCA for 2-D and 3-D visualization.

\begin{figure}
\centering
\begin{tabular}{@{}c@{\hspace{5mm}}c@{}}
\includegraphics[width=.4\linewidth]{ml/pca-nn/points-2d}&
\includegraphics[width=.4\linewidth]{ml/pca-nn/points-3d}
\end{tabular}
\caption{\label{fig-pca}
PCA projections of the data into two and three dimensions.
}
\end{figure}


   
\paragraph{Optimality analysis.}

\input{chapters/sec-pca-theory.tex}
 
 

                                                                  
\subsection{Clustering and \texorpdfstring{$k$}{k}-means}

Clustering assigns a label $y_i \in \{1,\ldots,k\}$ to each observation $x_i$. Observations with the same label form a cluster.

    
\paragraph{\texorpdfstring{$k$}{k}-means}

One approach is to seek compact clusters by minimizing a measure of within-cluster dispersion. In Euclidean space, the $k$-means objective measures distances from observations to their centroids $c=(c_\ell)_{\ell=1}^k$, with $c_\ell \in \RR^p$. Related formulations are possible in general metric spaces, although the centroid computations may be more difficult. The optimization problem is
\eq{
	\umin{ (y,c) } \Ee( y, c ) \eqdef \sum_{\ell=1}^k \sum_{ i : y_i=\ell } \norm{x_i-c_\ell}^2. 
}

\begin{figure}[tbp]
\centering
\BookDrawing{tikz/ml-voronoi}
\caption{$k$-means clusters according to Voronoi cells.}
\label{fig-review-ml-voronoi}
\end{figure}
The $k$-means algorithm alternates between updating the labels and updating the centroids, a form of block coordinate minimization.
 
Initialize the centroids $c$, for example by choosing well-separated observations. We discuss a more systematic initialization below.
  
For fixed centroids $c$, minimize $y \mapsto \Ee(y,c)$ by assigning each observation to its nearest centroid:
\eql{\label{eq-kmeans-1}
	\foralls i \in \{1,\ldots,n\}, \quad y_i \leftarrow \uargmin{1 \leq \ell \leq k} \norm{x_i-c_\ell}.
}
For fixed labels $y$, minimize $c \mapsto \Ee(y,c)$ by setting each centroid to the mean of its cluster:
\eql{\label{eq-kmeans-2}
	\foralls \ell \in \{1,\ldots,k\}, \quad c_\ell \leftarrow \frac{
		\sum_{i : y_i=\ell} x_i
	}{
		|\enscond{i}{y_i=\ell}|
	}
}
If a cluster becomes empty, its centroid does not contribute to the objective for the current labels. One may leave it in place or reseed it at a data point before the next assignment step. The 